% =====================================================================
%  Capítulo 3. Propuesta a desarrollar
% =====================================================================
\chapter[PROPUESTA A DESARROLLAR]{Propuesta a desarrollar}
\label{cap:propuesta}

Se propone el desarrollo de un sistema de comunicación asistida implementado en Python que permita
a personas con compromisos motores severos construir mensajes a partir de sus señales
electroencefalográficas, empleando el paradigma P300 Speller descrito en la
Sección~\ref{sec:speller}. El sistema presentará al usuario una matriz visual de caracteres cuyas
filas y columnas se iluminarán de manera aleatoria mientras este concentra su atención en el
carácter que desea seleccionar; las épocas de EEG asociadas a cada intensificación serán
clasificadas por un modelo de aprendizaje profundo como objetivo o no objetivo, y a partir de estas
clasificaciones se determinará el carácter elegido \cite{farwell1988}. Para agilizar la escritura,
el texto construido hasta ese momento alimentará un módulo de predicción de palabras que mostrará
sugerencias al usuario, de modo que una palabra completa pueda seleccionarse sin deletrearla letra
por letra, enfoque que ha mostrado reducir el tiempo de escritura en sistemas de este tipo
\cite{ryan2011,hong2026}.

Dado que la adquisición de señales propias implicaría contar con equipo especializado de 64
canales, los modelos se entrenarán y evaluarán con bases de datos públicas del paradigma P300,
principalmente la base de datos II de la tercera competencia internacional de BCI
\cite{blankertz2006}. Sobre estos datos se evaluarán estrategias para reducir el número de canales
y el número de repeticiones necesarias por carácter, con el propósito de acercar el sistema a
condiciones de uso más prácticas sin comprometer de manera importante su exactitud, que son las dos
limitaciones persistentes identificadas en el estado del arte.

\section[ARQUITECTURA DEL SISTEMA]{Arquitectura del sistema}

La arquitectura general del sistema se compone de las etapas que se muestran en la
Figura~\ref{fig:arquitectura}: la adquisición de las señales EEG a partir de bases de datos
públicas, su preprocesamiento, el entrenamiento del modelo de aprendizaje profundo, la
clasificación de los estímulos como objetivo o no objetivo y, finalmente, la construcción del texto
apoyada por el modelo de predicción de palabras, que mostrará sugerencias a partir de los caracteres
seleccionados hasta ese momento.

\begin{figure}[htbp]
\centering
\IfFileExists{imagenes/03-propuesta/arquitectura.png}{%
  \includegraphics[width=0.95\textwidth]{03-propuesta/arquitectura.png}%
}{%
  \fbox{\parbox[c][5cm][c]{0.9\textwidth}{\centering\small
  Colocar la imagen de la arquitectura del sistema en\\
  \texttt{imagenes/03-propuesta/arquitectura.png}}}%
}
\caption{Arquitectura del sistema. Elaboración propia.}
\label{fig:arquitectura}
\end{figure}

La etapa de preprocesamiento aplicará las técnicas descritas en la Sección~\ref{sec:pds}, entre
ellas el filtrado pasa banda, el remuestreo, la segmentación en ventanas temporales posteriores a
cada estímulo y la normalización de cada canal, con el fin de generar un conjunto de datos adecuado
para el entrenamiento y la validación del modelo. El modelo de detección se construirá con la
biblioteca PyTorch, y para el procesamiento y manipulación de los datos se emplearán bibliotecas
como NumPy y Pandas. La interfaz gráfica, por su parte, se desarrollará con PyQt e integrará la
matriz visual del P300 Speller, el área en la que se muestra el texto construido y las sugerencias
del módulo de predicción de palabras.

El sistema se implementará inicialmente en una computadora con sistema operativo Windows, en la
que se realizará el entrenamiento, la validación y la inferencia de los modelos. Así mismo, se
analizará la viabilidad de ejecutar únicamente la etapa de inferencia en un microordenador NVIDIA
Jetson, con el propósito de evaluar la posibilidad de una implementación más portable, en línea con
el uso de arquitecturas compactas como EEGNet o SepConv1D, cuyo reducido número de parámetros las
hace convenientes para dispositivos con recursos limitados \cite{lawhern2018,alvarado2021}.

\section[FLUJO DE OPERACIÓN]{Flujo de operación}

Durante el uso del sistema, la selección de cada carácter sigue el ciclo que se representa en la
Figura~\ref{fig:flujo}. Las filas y columnas de la matriz se intensifican en orden aleatorio y el
usuario mantiene su atención en el carácter que desea escribir; cuando la fila o la columna que lo
contiene se ilumina, se genera un potencial P300 aproximadamente 300~ms después del estímulo
\cite{farwell1988,polich2007}. La señal registrada se divide en épocas sincronizadas con cada
intensificación, que después de ser preprocesadas se entregan al modelo, el cual asigna a cada una
un puntaje que refleja la probabilidad de que contenga una P300. Estos puntajes se acumulan por
fila y por columna a lo largo de las repeticiones, conforme a la Ecuación~\eqref{eq:decodificacion},
y el carácter seleccionado corresponde a la intersección de la fila y la columna con el puntaje más
alto.

\begin{figure}[htbp]
\centering
\begin{tikzpicture}[
  node distance=0.55cm and 0.5cm,
  paso/.style={draw, rounded corners=2pt, fill=gray!8, align=center, minimum height=1.15cm,
               text width=2.55cm, font=\footnotesize},
  flecha/.style={-{Stealth[length=2mm]}, thick}]
  \node[paso, fill=red!8] (est) {Intensificación aleatoria de filas y columnas};
  \node[paso, right=of est] (eeg) {Registro del EEG y segmentación en épocas};
  \node[paso, right=of eeg] (pre) {Preprocesamiento (filtrado, remuestreo, normalización)};
  \node[paso, right=of pre, fill=blue!8] (mod) {Modelo de aprendizaje profundo: puntaje por época};
  \node[paso, below=0.9cm of mod] (dec) {Acumulación de puntajes y selección del carácter};
  \node[paso, left=of dec, fill=blue!8] (pal) {Predicción de palabras a partir del texto escrito};
  \node[paso, left=of pal] (sug) {Presentación de sugerencias en la interfaz};
  \node[paso, left=of sug, fill=red!8] (usr) {Usuario elige un carácter o una palabra sugerida};
  \draw[flecha] (est) -- (eeg);
  \draw[flecha] (eeg) -- (pre);
  \draw[flecha] (pre) -- (mod);
  \draw[flecha] (mod) -- (dec);
  \draw[flecha] (dec) -- (pal);
  \draw[flecha] (pal) -- (sug);
  \draw[flecha] (sug) -- (usr);
  \draw[flecha, dashed] (usr) -- node[left, font=\scriptsize, align=right] {siguiente\\selección} (est);
\end{tikzpicture}
\caption{Flujo de operación para la selección de un carácter o de una palabra sugerida.
Elaboración propia.}
\label{fig:flujo}
\end{figure}

Una vez que el carácter se agrega al texto, el módulo de predicción de palabras estima las
palabras más probables a partir de los caracteres escritos, siguiendo alguno de los enfoques de
modelado del lenguaje revisados en la Sección~\ref{sec:lenguaje}, y las presenta en la interfaz
como opciones adicionales que el usuario puede seleccionar mediante el mismo mecanismo de
atención. De esta forma, cada selección puede corresponder a un carácter o a una palabra completa,
y el número de repeticiones empleado en cada una de ellas será uno de los parámetros a optimizar,
pues determina directamente la tasa de transferencia de información del sistema, tal como se
describió en la Sección~\ref{sec:metricas}.

\section[METODOLOGÍA]{Metodología}

El desarrollo del presente Trabajo Terminal se realizará mediante una combinación de dos
metodologías, dado que el proyecto involucra tanto la construcción de un sistema de software como
la de un modelo de aprendizaje profundo. Para el desarrollo general del sistema se empleará un
modelo de desarrollo incremental, mientras que para el modelo de detección de la componente P300 se
utilizará la metodología CRISP-ML, la cual estructura el proceso de construcción, validación y
optimización de modelos de aprendizaje automático en fases iterativas \cite{studer2021}. Ambas
metodologías se ejecutarán en paralelo y convergerán en la etapa de integración, en la que el
modelo entrenado se incorpora al sistema completo.

\subsection{Metodología de desarrollo del sistema}

En el modelo incremental, el software se construye de manera progresiva mediante una serie de
incrementos, cada uno de los cuales agrega funcionalidad a los anteriores y produce una versión
operativa que puede evaluarse antes de continuar con el siguiente \cite{pressman2010}. En este
trabajo, cada incremento corresponderá a la implementación e integración de uno de los módulos
funcionales del sistema, lo que permitirá evaluar sus componentes de forma progresiva mientras el
modelo de detección de la componente P300 se desarrolla en paralelo. La
Figura~\ref{fig:incremental} ilustra la organización de los incrementos, y las etapas consideradas
son las siguientes:

\begin{enumerate}
  \item \textbf{Definición de requerimientos del sistema.} En esta etapa se realizará una revisión
  del estado del arte y se identificarán los requerimientos funcionales y no funcionales del
  sistema de comunicación asistida, así como las características que deberá cumplir la interfaz de
  usuario y los módulos que integrarán la aplicación.
  \item \textbf{Diseño de la arquitectura del sistema.} Se definirá la arquitectura general,
  identificando los módulos que la conforman y las interacciones entre ellos; entre estos se
  consideran el módulo de detección de la componente P300, el módulo de predicción de palabras y
  la interfaz de usuario.
  \item \textbf{Diseño de la interfaz de usuario.} Se desarrollarán los prototipos de la interfaz
  gráfica que permitirá la interacción del usuario con la matriz del P300 Speller y con las
  sugerencias del módulo de predicción de palabras.
  \item \textbf{Desarrollo e integración de módulos.} Se implementarán los distintos módulos del
  sistema, incluyendo la interfaz de usuario, el módulo de predicción de palabras y la integración
  del modelo de detección de la componente P300, agregando cada uno como un nuevo incremento sobre
  la versión anterior.
  \item \textbf{Pruebas funcionales del sistema.} Se verificará el funcionamiento de la
  interacción entre los módulos, así como la correcta construcción de palabras mediante el
  paradigma P300 Speller.
  \item \textbf{Documentación del sistema.} Finalmente, se generará la documentación técnica del
  sistema y el manual de usuario correspondiente.
\end{enumerate}

\begin{figure}[htbp]
\centering
\begin{tikzpicture}[
  fase/.style={draw, minimum height=0.7cm, minimum width=1.75cm, font=\scriptsize, align=center},
  flecha/.style={-{Stealth[length=1.8mm]}, thick}]
  \foreach \i/\nombre in {1/{Incremento 1: interfaz y matriz P300}, 2/{Incremento 2: predicción de palabras}, 3/{Incremento 3: integración del modelo P300}}{
    \pgfmathsetmacro{\x}{(\i-1)*2.2}
    \pgfmathsetmacro{\y}{-(\i-1)*1.35}
    \node[fase, fill=gray!10, anchor=west] (a\i) at (\x,\y) {Análisis};
    \node[fase, fill=blue!8, right=0cm of a\i] (d\i) {Diseño};
    \node[fase, fill=red!8, right=0cm of d\i] (c\i) {Codificación};
    \node[fase, fill=green!10, right=0cm of c\i] (p\i) {Pruebas};
    \node[font=\scriptsize, anchor=south west] at (a\i.north west) {\nombre};
    \node[font=\scriptsize, right=0.08cm of p\i] {Versión \i};
  }
  \draw[flecha] (0,-3.4) -- node[below, font=\scriptsize] {Tiempo de desarrollo} (13.4,-3.4);
\end{tikzpicture}
\caption{Organización del desarrollo del sistema mediante el modelo incremental. Elaboración propia
con base en \cite{pressman2010}.}
\label{fig:incremental}
\end{figure}

\subsection{Metodología para el desarrollo del modelo de aprendizaje profundo}

El modelo de detección de la componente P300 se desarrollará siguiendo la metodología CRISP-ML,
una adaptación del proceso CRISP-DM para proyectos de aprendizaje automático que organiza el
trabajo en fases con retroalimentación entre ellas, de modo que es posible regresar a la
preparación de los datos o al modelado cuando los resultados de la evaluación no son
satisfactorios \cite{studer2021}. La Figura~\ref{fig:crispml} muestra el ciclo de estas fases, que
en el presente trabajo se describen a continuación:

\begin{enumerate}
  \item \textbf{Comprensión del problema.} Se establecerán los objetivos particulares de detección
  del modelo y los criterios con los que se evaluará su desempeño.
  \item \textbf{Comprensión y análisis de los datos.} Se seleccionarán bases de datos públicas de
  señales EEG asociadas al paradigma P300 Speller, analizando su estructura, número de sujetos y
  cantidad de canales, con el fin de determinar las características necesarias para el
  entrenamiento del modelo.
  \item \textbf{Preparación de los datos.} Se aplicarán técnicas de preprocesamiento como el
  filtrado pasa banda, la división en ventanas temporales y la normalización, generando un
  conjunto de datos adecuado para el entrenamiento y la validación del modelo de detección de la
  componente P300.
  \item \textbf{Diseño y modelado.} Se definirán las posibles arquitecturas de aprendizaje profundo
  para la detección de la componente P300, se realizará el ajuste de hiperparámetros y se
  llevarán a cabo experimentos orientados a la reducción del número de canales y de repeticiones
  necesarias para la detección.
  \item \textbf{Evaluación.} Se analizará el desempeño del modelo mediante métricas de
  clasificación como la exactitud, la sensibilidad, la precisión y el F1-score, además de la
  construcción de la matriz de confusión, y se evaluará el impacto de la reducción de canales y
  repeticiones en el rendimiento del sistema.
  \item \textbf{Integración del modelo.} El modelo desarrollado se integrará en el sistema completo
  de comunicación asistida para su uso dentro del paradigma P300 Speller.
\end{enumerate}

\begin{figure}[htbp]
\centering
\begin{tikzpicture}[
  fase/.style={draw, rounded corners=3pt, fill=blue!6, align=center, text width=2.4cm,
               minimum height=1.0cm, font=\footnotesize},
  flecha/.style={-{Stealth[length=2mm]}, thick}]
  \def\r{3.3}
  \foreach \i/\ang/\texto in {1/90/Comprensión del problema, 2/30/Comprensión de los datos,
      3/-30/Preparación de los datos, 4/-90/Diseño y modelado, 5/-150/Evaluación,
      6/150/Integración del modelo}{
    \node[fase] (f\i) at (\ang:\r) {\texto};
  }
  \node[circle, draw, fill=red!8, align=center, font=\footnotesize, minimum size=2.1cm] (c)
    at (0,0) {Datos EEG\\públicos};
  \foreach \a/\b in {1/2, 2/3, 3/4, 4/5, 5/6}{
    \draw[flecha] (f\a) to[bend left=12] (f\b);
  }
  \draw[flecha, dashed] (f5.east) -- node[below, font=\scriptsize] {reajuste} (f3.west);
  \draw[flecha, dashed] (c) -- (f2);
\end{tikzpicture}
\caption{Fases de la metodología CRISP-ML aplicadas al desarrollo del modelo de detección de la
componente P300. Elaboración propia con base en \cite{studer2021}.}
\label{fig:crispml}
\end{figure}

Debido a que el proyecto requiere una integración estrecha entre el procesamiento de las señales
EEG, el desarrollo del modelo de aprendizaje profundo y la implementación del sistema de
comunicación asistida, algunas actividades relacionadas con el diseño experimental, el desarrollo y
la evaluación del modelo serán realizadas por ambos integrantes. No obstante, cada integrante
tendrá un mayor enfoque en determinadas áreas: uno se orientará principalmente al desarrollo de la
interfaz gráfica, la integración de módulos y la arquitectura general del sistema, mientras que el
otro se enfocará en el análisis de los datos EEG, la experimentación y la evaluación del desempeño
del modelo.
